\chapter{Conclusion}
\label{ch:conclusion-en}

This conclusion introduces no new empirical claims. It closes only chains already
assigned statuses in C01--C11 and preserves the distinction among
\code{supported}, \code{rejected}, \code{descriptive}, and
\code{not_tested}.

The dissertation aimed to compare the investigated predictive-coding regimes with BP
within a frozen experimental domain, move from final metrics to internal computational
organization, and test whether state available before normal completion could support
an early compute-control decision under a registered dangerous-accept criterion and
positive economic benefit. The work therefore proceeded through Stage~1/2, Stage~3A,
Stage~3B, and QWake-FP C1/C2 with versioned protocols, separate admission gates, and
preserved artifact provenance.

\section{The PC-TREF, PC-CATM, and QWake-PC Theoretical Connection}

Alongside its empirical results, the dissertation formulates a connection among
separate theoretical levels linking the internal PC mechanism to a compute-budget
decision. PC-CATM decomposes the resulting correction into structurally distinct
channels, transport, and cancellation. PC-TREF uses a diagnostic representation of
states and distinguishes the strong criterion of full representation sufficiency on
the registered domain from narrower selective sufficiency for a specific action.
QWake-PC occupies the next level and turns these criteria and available features into
a proposal for a computational action.

In QWake-FP, the required partition distinguishes \code{EARLY_ADMISSIBLE} from
\code{CONTINUE_REQUIRED}. The first class means that the required response of the
registered \code{ANALYTIC_COMPLETION}
\nolinkurl{fixedpred_eta1_wavefront_completion_v1} is equivalent, within the
registered tolerance, to the required response of the exact suffix
\nolinkurl{complete_suffix_stage2_baseline_v1}. The early action therefore replaces
the canonical iterative suffix with another computational path; it does not remove all
further computation.

These levels are intentionally not identified with one another. A PC-CATM diagnostic
regime is not permission to take an early action; PC-TREF sufficiency does not guarantee
economic benefit; and QWake-PC is not one fixed algorithm. The C1/C2 empirical result
therefore concerns QWake-FP, a particular bounded shadow instantiation for FixedPred,
and does not automatically transfer to the whole architecture or other PC regimes.
PC-CATM remains theoretical/diagnostic in the current work: comparative superiority of
its NCZ/ECZ/TNZ or other mechanistic features for compute control was not directly
tested.

\section{Answer to RQ1}
\label{sec:conclusion-rq1-en}

RQ1 asked under which frozen conditions Exact, FixedPred, and Strict are close to BP
and where differences exceed selected boundaries.

Stage~1/2 shows that similarity of final quality can coexist with differences in
computational cost. In the registered paired FashionMNIST analysis, FixedPred and
Strict fall inside the preregistered test-macro-F1 equivalence margin relative to BP.
Stage~2 preserves the registered final-quality surface while changing runtime relative
to BP.

Stage~3A bounds the interpretation of this behavioral proximity. In the registered
domain, FixedPred shows higher observed similarity to BP than Strict in gradient
direction and representation metrics, while early-gradient norm is substantially
reduced. The RQ1 answer is therefore not binary ``equivalent/not equivalent''.
Proximity depends on the level of analysis---final metric, gradient direction, gradient
magnitude, or representation state---and must be reported with its corresponding
inference boundary.

C01 records the supported Stage~1 paired macro-F1 equivalence result together with the
Stage~2 final-quality preservation and cross-version runtime change; C02 independently
records the supported layer-wise Stage~3A result.

\section{Answer to RQ2}
\label{sec:conclusion-rq2-en}

RQ2 concerned localization of computational differences and their relation to
mechanisms and cost components.

Stage~3B B0 localizes \texttt{state\_inference} as the principal device-time region in
the studied ROCm/float32 environment; Strict has higher device time and peak memory
than FixedPred on this matrix. B0 uses three independent \texttt{model\_seed} values
per configuration, so C03 is interpreted as descriptive localization on the
registered profiling surface rather than a precise population estimate.

The primary SI-MA0 mechanistic cost-attribution model fails COST-MA0 and that result is
preserved as C04. Independent SI-MA1 calibration passes CAL-COST-MA1; its signed
residual is interpreted as accounting over-closure, not negative physical cost,
supporting C05.

Exact B1 and B2 implementation candidates pass their own equivalence gates before
comparative profiling, supporting C06. The registered matched 288-cell matrix then
completes without failed or repeated runs. However, the separate registered
engineering-continuation screen qualifies no configuration in any of the four
B1/B2-by-FixedPred/Strict groups: each yields 0/16 and both candidates receive
\texttt{reject\_or\_revise}. C07 records this bounded negative result descriptively.
Passing numerical and trajectory equivalence does not guarantee passage of the
resource-continuation screen, but neither universal baseline superiority nor universal
non-viability of B1/B2 is inferred.

RQ2 therefore receives a sequential-localization answer: cost is associated with the
specific \texttt{state\_inference} region; the initial attribution hypothesis receives
a negative result; observer-cost calibration is tested separately; and alternative
exact computational organizations reach matched profiling after functional-equivalence
gates but receive \texttt{reject\_or\_revise} on the registered engineering screen.

\section{Answer to RQ3}
\label{sec:conclusion-rq3-en}

RQ3 combined three sequential levels: existence of preterminal records for which the
registered analytic completion is admissible relative to the required response of the
exact suffix; ability to recognize a non-zero subset on the frozen calibration surface
from pre-action state with zero observed dangerous accepts; and positive computational
benefit.

On the sealed QWake-FP calibration surface, the first two levels hold in a bounded
selective sense. Structural reconciliation gives 108 trajectories with seven epochs
\(t=0,\ldots,6\). The highest-coverage zero-danger rule,
\texttt{compute\_step >= 5}, accepts 108 preterminal step-5 records with one remaining
sweep and 108 terminal-boundary step-6 records; all 216 accepted records have zero
registered dangerous accepts. C08 therefore rests on at least 108 explicitly
preterminal accepted records whose analytic candidate is labeled
\code{EARLY_ADMISSIBLE} by the oracle. But the rule itself uses only
\texttt{compute\_step >= 5}; C08 establishes a temporal fixed-prefix boundary on this
surface and does not demonstrate input-dependent adaptivity or superiority of
PC-CATM features. Registered 216/756 coverage (about 28.57\%) concerns the complete
calibration surface and is not preterminal early-action coverage. C08 is therefore
\texttt{supported} only within this calibration boundary; observing zero dangerous
accepts does not establish population-level safety.

The third, economic condition fails under frozen full decision-cost accounting. None
of the 2,625 rules has zero dangerous accepts, non-zero coverage, and positive
aggregate net saving simultaneously. C09 is \texttt{rejected}. For the highest-
coverage zero-danger rule, full decision cost greatly exceeds the registered
\texttt{remaining\_suffix\_ns} sum over accepted non-dangerous records, and about
99.887\% of the accounted cost belongs to the observer category. The original C2
estimand is retained; the residual-suffix sum is not interpreted as pure avoidable
sweep time because accepted records include the terminal boundary.

This negative economic result does not transfer to marginal execution cost of a simple
recognizer. That estimand was outside C2, so C10 remains \code{not_tested}. Because
no eligible rule existed, C2 established no rule freeze and the protocol did not open
C3; C11 also remains \code{not_tested}.

RQ3 therefore receives a three-part answer: preterminal states admitting the
registered early action exist on the studied surface; 108 such step-5 records lie in
the subset accepted by the highest-coverage zero-danger rule, while the original
216/756 registered coverage additionally includes 108 terminal-boundary records; but
no zero-danger rule with positive aggregate net saving exists under the frozen full
accounting. Economic viability of a minimal implementation was not tested.

\section{Scientific and Methodological Outcome}
\label{sec:conclusion-contribution-en}

Within its declared domain, the work links analytical levels often treated separately:
final behavior, internal gradients and representations, mechanistic profiling, and a
bounded decision about residual computation. The principal result is not selection of
one ``best'' regime but localization of which properties agree, which differ, and which
next hypothesis becomes admissible after each stage.

Methodologically, the study separates difference from equivalence claims, uses the
independently trained model as statistical unit, preserves negative tests, separates
pre-action from post-action information in QWake-FP, and gates sequential stages by
frozen receipts. Thesis-facing summaries bind key numerical claims to tracked evidence
through SHA-256 and are checked in CI before LaTeX generation.

This matters especially for negative results. SI-MA0 and C2 do not disappear after
narrower explanations are found. They remain constraints on their respective
hypotheses and determine whether a next experiment is scientifically new or merely a
post-hoc modification of a completed estimand.

\section{Inference Boundaries}

The results concern the investigated Torch2PC versions, datasets,
\code{lenet\_classic}, model seeds, algorithmic regimes, numerical tolerances, and
frozen hardware/software environments. They do not establish universal PC/BP
equivalence or a universally preferred implementation.

QWake-FP C2 is additionally bounded to a finite rule family, 756 sealed calibration
records, and frozen full decision-cost accounting. Zero dangerous accepts on this
surface are not a population safety guarantee. Failure to find a zero-danger,
economically beneficial rule does not exclude another recognizer, another state
representation, or another architecture with lower marginal cost.

These boundaries are not post-hoc caveats: the statuses \code{supported},
\code{rejected}, \code{descriptive}, and \code{not_tested} are part of the
claim contract and constrain the wording of conclusions.

\section{Future Work}

The immediate next research question is localized by C08--C10: can the observed or a
functionally comparable recognizer preserve a preregistered dangerous-accept criterion
on a new confirmatory surface at sufficiently low marginal cost that its actual extra
runtime is below the cost of the avoided residual computation? This requires a new
preregistration, a new cost-measurement surface, and independent measurement; the old
C2 must not be retrospectively recomputed under a new accounting model.

Once an eligible rule with measured marginal cost exists, an independent confirmatory
campaign can be run on a new evidence surface. It requires its own protocol identity
and is not a retroactively opened C3 of the original C2 chain.

Further work includes transfer of the sequential analysis to other architectures,
datasets, numerical-precision regimes, and hardware, and investigation of cheaper
observer representations. In every case, behavioral similarity, internal mechanism,
dangerous-accept criterion, and economic cost should remain separate.

\section{Final Conclusion}

The completed experimental program progressively narrowed the original question about
the relation between predictive coding and backpropagation. In the studied domain,
proximity of final quality does not imply identity of gradients, representations, or
runtime; profiling localizes substantial cost in \texttt{state\_inference}; B1/B2 show
that numerical and trajectory admissibility do not guarantee passage of a separate
resource-continuation screen; and QWake-FP shows that a pre-action signal selecting a
non-zero calibration region with zero observed dangerous accepts does not by itself
establish economic viability.

The separation of these levels is the principal outcome: bounded informational
viability is supported while economic viability under the frozen accounting is
rejected, and economic viability of a minimal implementation remains an open but now
precisely formulated next question. The completed evidence chain reaches its
preregistered stopping boundary: C2 remains bounded-negative on the economic test, C3
remains closed, and any new claim requires a new experimental contract.
